% Part: lambda-calculus
% Chapter: lambda-definability
% Section: fixpoints

\documentclass[../../../include/open-logic-section]{subfiles}

\begin{document}

\olfileid{lam}{ldf}{fp}

\olsection{مقرر نقطے}

فرض کرو اسیں عاملی فنکشن بازگشت نال ٹرم $\fn{Fac}$ دے طور
اُتے متعین کرنا چاہندے آں، جیہدی خاصیت ایہ ہووے۔
\[
\fn{Fac} \ident \lambd[n][\fn{IsZero}\, n\, \num 1 (\fn{Mult}\, n (\fn{Fac}(\fn{Pred} \, n)))]
\]
یعنی $n$ دا عامل $1$ اے جے $n = 0$، تے ہور صورت وچ
$n$ گُنا $n-1$ دا عامل۔ صاف اے کہ ٹرم $\fn{Fac}$ نوں
ایویں متعین نہیں کر سکدے، کیوں جے سجے پاسے $\fn{Fac}$ آپ
وی آؤندا اے۔ اپنے آپ دا حوالہ دین والیاں ایہو جیہیاں بازگشتی
تعریفاں لامبڈا کیلکولس دا حصہ نہیں۔ مثلاً ٹرم دی ایس تعریف وچ
\[
\fn{Mult} \ident \lambd[ab][a (\fn{Add}\, a) 0]
\]
سجے پاسے صرف پہلاں توں متعین ٹرم آؤندے نیں، جیویں $\fn{Add}$۔
$\fn{Add}$ نوں اوہدی تعریف والا ٹرم رکھ کے ہمیشہ ہٹا سکدے آں۔
ایس نال $\fn{Mult}$ خالص لامبڈا ٹرم بن جاندا اے۔ جے $\fn{Add}$
وچ وی متعین ٹرم ہون، جیویں $\fn{Add}'$ وچ نیں، تاں ایس عمل نوں
جاری رکھ کے آخر اک خالص لامبڈا ٹرم تک پہنچ سکدے آں۔

پر اُپر دِتّی $\fn{Fac}$ ورگی بازگشتی تعریف لئی ایہ گل
سچ نہیں۔ جے سجے پاسے $\fn{Fac}$ دی آمد دی تھاں $\fn{Fac}$
دی اپنی تعریف رکھیے، تاں سانوں ایہ ملدا اے۔
\begin{align*}
  \fn{Fac} & \ident \lambd[n][\fn{IsZero}\, n\, \num{1}] \\
    & \qquad (\fn{Mult}\, n
  ((\lambd[n][\fn{IsZero} \, n \, \num 1\, (\fn{Mult}\, n\, (\fn{Fac}
    (\fn{Pred}\, n)))]) (\fn{Pred}\, n)))
\end{align*}
تے سجے پاسے $\fn{Fac}$ ہلے وی موجود اے۔ صاف اے کہ ایہ عمل
دہراؤن نال تعریف لمّی ہوندی جاندی اے تے کدے خالص لامبڈا ٹرم
نہیں ملدا۔ سو عاملی فنکشن، یا عام طور اُتے بازگشتی فنکشن،
ایس طریقے نال متعین نہیں کیتے جا سکدے۔

فیر وی بازگشتی تعریف توں اک گل پتا لگدی اے۔ جے $f$ عاملی
فنکشن دا ٹرم ہووے، تاں ہِٹھلا ٹرم
\[
\fn{Fac}' \ident \lambd[g][\lambd[n][\fn{IsZero} \, n \, \num 1\, (\fn{Mult} \, n\, (g (\fn{Pred} n)))]]
\]
ٹرم $f$ اُتے لاؤن نال، یعنی $\fn{Fac}'\,f$، عاملی فنکشن ای
ظاہر کردا اے۔ یعنی جے $\fn{Fac}'$ نوں ایہو جیہا فنکشن
سمجھیے جیہڑا فنکشن لیندا تے فنکشن واپس دیندا اے، تاں
$\fn{Fac'}\, f$ دی قدر بس~$f$ اے، بشرطیکہ $f$ عاملی
فنکشن ہووے۔ فنکشن~$f$ نوں $\fn{Fac}' \, f \equal[\beta] f$
ہوون اُتے $\fn{Fac}'$ دا \emph{مقرر نقطہ} آکھدے آں۔
سو عاملی فنکشن $\fn{Fac}'$ دا مقرر نقطہ اے۔

لامبڈا کیلکولس وچ ایہو جیہے ٹرم نیں جیہڑے دِتّے ٹرم دے مقرر
نقطے کڈھدے نیں۔ ایہناں نال $\fn{Fac}'$ ورگے ٹرم توں عاملی
فنکشن دی تعریف بنائی جا سکدی اے۔

\begin{defn}
  \ollabel{defn:Turing-Y}
  \emph{وائی کمبینیٹر} ایہ ٹرم اے۔
  \[
  Y \ident (\lambd[ux][x(uux)])(\lambd[ux][x(uux)]).
  \]
\end{defn}

\begin{thm}
  $Y$ دی خاصیت اے کہ $Yg \red g(Yg)$ ہر ٹرم $g$ لئی اے۔
  سو $Yg$ ہمیشہ~$g$ دا مقرر نقطہ اے۔
\end{thm}

\begin{proof}
  $(\lambd[ux][x(uux)])$ نوں مختصر~$U$ لکھ لئیے، تاں جو
  $Y \ident UU$۔ فیر
  \begin{align*}
    Y g &\ident (\lambd[ux][x(uux)])U\,g \\
    &\red (\lambd[x][x(UUx)])g\\
    & \red g(UUg) \ident g(Yg).
  \end{align*}
  کیوں جے $g(Yg)$ تے $Yg$ دوہاں دی تقلیل $g(Yg)$ تک ہوندی اے،
  ایس لئی $g(Yg) \equal[\beta]
  Yg$، تے $Yg$،~$g$ دا مقرر نقطہ اے۔
\end{proof}

کیوں جے $Yg$ اک ریڈیکس اے، ایس لئی تقلیل بغیر مُکے جاری رہ
سکدی اے۔
\begin{align*}
  Y g &\red g (Y g) \\
    &\red g (g (Y g)) \\
    &\red g(g (g (Y g)))\\
    &\ldots
\end{align*}
سو اسیں $Yg$ نوں $g$ دا اپنے آپ اُتے لامتناہی وار اطلاق
سمجھ سکدے آں۔ جے~$g$ نوں ایس اُتے اک وار ہور لائیے، تاں
ایویں آکھ لو کہ کوئی وادھا نہیں ہوندا۔ $g$ نوں~$g$ دے
~$Yg$ اُتے لامتناہی وار اطلاق دے نتیجے اُتے لاؤنا ہلے وی
$g$ دا~$Yg$ اُتے لامتناہی وار اطلاق ای اے۔

دھیان دیو کہ $\beta$-تقلیل دے قدماں دا اُپرلا سلسلہ
~$Yg$ توں شروع ہوندا اے تے لامتناہی اے۔ سو جے $Yg$ نوں
کسے ٹرم اُتے لائیے، یعنی $(Yg)N$ ویکھیے، تاں اوہ ٹرم وی
لامتناہی وکھرے ٹرماں تک گھٹدا اے، یعنی $(g(Yg))N$،
$(g(g(Yg)))N$، \dots۔ فیر وی ہو سکدا اے کہ تقلیل دے
قدماں دا کوئی \emph{ہور} سلسلہ نارمل روپ اُتے مُک جاوے۔

مثال لئی عاملی فنکشن ویکھو۔ $\fn{Fac}$ دی تعریف
$Y\, \fn{Fac}'$ دے طور اُتے کرو، یعنی $\fn{Fac'}$ دے مقرر
نقطے وجوں۔ فیر
\begin{align*}
  \fn{Fac}\,\num 3
  &\red Y\, \fn{Fac}' \, \num 3 \\
  &\red \fn{Fac}' (Y \,\fn{Fac}') \, \num 3 \\
  &\ident (\lambd[x][\lambd[n][\fn{IsZero} \, n\, \num 1\,
      (\fn{Mult}\, n\, (x (\fn{Pred}\, n)))]]) \, \fn{Fac} \, \num 3\\
  & \red \fn{IsZero} \, \num 3\, \num 1\,
      (\fn{Mult}\, \num 3\, (\fn{Fac} (\fn{Pred}\, \num 3))) \\
  &\red  \fn{Mult}\, \num 3 \, (\fn{Fac} \, \num 2).
  \intertext{ایسے طرح}
  \fn{Fac}\,\num 2
  &\red  \fn{Mult}\, \num 2 \, (\fn{Fac} \, \num 1) \\
  \fn{Fac}\,\num 1
  &\red  \fn{Mult}\, \num 1 \, (\fn{Fac} \, \num 0)
  \intertext{پر}
  \fn{Fac}\,\num 0
  &\red \fn{Fac}' (Y \,\fn{Fac}') \, \num 0 \\
  &\ident (\lambd[x][\lambd[n][\fn{IsZero} \, n\, \num 1\,
      (\fn{Mult}\, n\, (x (\fn{Pred}\, n)))]]) \, \fn{Fac} \, \num 0\\
  & \red \fn{IsZero} \, \num 0\, \num 1\,
      (\fn{Mult}\, \num 0\, (\fn{Fac} (\fn{Pred}\, \num 0))).\\
  &\red \num 1.
  \intertext{سو سارے قدم جوڑ کے}
  \fn{Fac}\,\num 3
  &\red  \fn{Mult}\, \num 3 \,
  (\fn{Mult} \, \num 2\, (\fn{Mult} \, \num 1\, \num 1)).
\end{align*}

جیہڑی گل $\fn{Fac'}$ لئی اے، اوہی کسے وی بازگشتی تعریف لئی
اے۔ فرض کرو ساڈے کول اک بازگشتی مساوات اے۔
\begin{align*}
g\,x_1\dots x_n & \equal[\beta] N
\intertext{ایتھے $N$ وچ~$g$ تے $x_1$، \dots،~$x_n$ آ سکدے نیں۔
  فیر ہمیشہ اک ایہو جیہا ٹرم~$G \ident (Y \lambd[g][\lambd[x_1\dots x_n][N]])$ اے جیہدے لئی}
G\,x_1 \dots x_n & \equal[\beta] \Subst{N}{G}{g}.
\intertext{کیوں جے مقرر نقطے دے قضیے مطابق}
G \ident (Y \lambd[g][\lambd[x_1\dots x_n][N]]) & \red \lambd[g][\lambd[x_1\dots x_n][N]](Y \lambd[g][\lambd[x_1\dots x_n][N]])\\
& \ident (\lambd[g][\lambd[x_1\dots x_n][N]])\,G
\intertext{تے ایس توں}
G\,x_1 \dots x_n
& \red (\lambd[g][\lambd[x_1\dots x_n][N]])\,G\,x_1\dots x_n\\
& \red (\lambd[x_1\dots x_n][\Subst{N}{G}{g}])\,x_1\dots x_n\\
  & \red \Subst{N}{G}{g}.
\end{align*}

$Y$ کمبینیٹر، جیہڑا \olref{defn:Turing-Y} وچ دِتّا اے، ایلن ٹیورنگ دا اے۔ الونزو
چرچ نے اک ہور روپ تجویز کیتا سی جیہنوں اسیں~$Y_C$ آکھاں گے۔
\[
Y_C \ident \lambd[g][(\lambd[x][g(xx)])(\lambd[x][g(xx)])].
\]
چرچ دا کمبینیٹر ٹیورنگ دے کمبینیٹر نالوں ذرا کمزور اے۔
ایہدے لئی $Yg
\equal[\beta] g(Yg)$ اے، پر $Yg \bred g(Yg)$ نہیں۔ منّ لو $V$
ٹرم $\lambd[x][g(xx)]$ اے، تاں جو $Y_C \ident \lambd[g][VV]$۔
فیر
\begin{align*}
  VV & \ident (\lambd[x][g(xx)])V \red g(VV)
  \text{ تے ایس لئی}\\
  Y_C g & \ident (\lambd[g][VV])g \red VV \red g(VV), \text{ پر ایہ وی}\\
  g(Y_C g) & \ident g((\lambd[g][VV])g) \red g(VV).
\end{align*}
دوجے لفظاں وچ، $Y_Cg$ تے $g(Y_Cg)$ دوہاں دی تقلیل اکّو ٹرم
$g(VV)$ تک ہوندی اے۔ سو $Y_Cg
\equal[\beta] g(Y_Cg)$۔ کئی ورتوں لئی ایہ کافی ہوندا اے۔

\end{document}
